% Responsável: TODO(grupo)
% Enunciado, item 6.3: no MÍNIMO 2 referências bibliográficas POR TÉCNICA.
% Rubrica: "Levantamento e seleção dos métodos" (peso 1).
\chapter{Fundamentação teórica}
\label{cap:fundamentacao}

\section{Assinatura vibratória de defeitos em rolamentos}
\label{sec:assinatura}

TODO(grupo): impacto periódico, modulação da ressonância estrutural e por que o
espectro bruto costuma não revelar a falha.

\section{Frequências características do rolamento}
\label{sec:frequencias-caracteristicas}

TODO(grupo): BPFO, BPFI, BSF e FTF a partir da geometria do 6205 e da rotação do
eixo. Este é o argumento físico mais forte disponível para o relatório --- mostrar
que o pico do envelope cai na frequência prevista pela teoria.

\section{Técnicas selecionadas}
\label{sec:tecnicas}

% Uma subseção por técnica. Cada uma precisa de 2 referências próprias.
\subsection{Análise de envelope (demodulação por Hilbert)}
TODO(grupo). Referências: TODO, TODO.

\subsection{Detecção one-class}
TODO(grupo). Referências: TODO, TODO.

\subsection{Classificação supervisionada}
TODO(grupo). Referências: TODO, TODO.
